\section{Experiments}
\label{sec:exp}

We answer six research questions.
\textbf{RQ1}: How many classification dimensions does the corpus support, how
widely is each shared across sources, and how faithful are the citations?
\textbf{RQ2}: What hierarchy do the term types form, how confident and how
grounded are its edges, and how does it change as weak edges are hidden?
\textbf{RQ3}: Does the DAG serve as a label space for Web content, and does
the content reveal where it must grow?
\textbf{RQ4}: How does the induced ontology compare with expert-built
taxonomies and with LLM baselines that are given the same corpus?
\textbf{RQ5}: How much do the manifest, the grounding strategy and the model
matter?
\textbf{RQ6}: Does the ontology improve a downstream classification of
Web-sourced token descriptions?
RQ1 to RQ3 and the manifest part of RQ5 are answered below; the remaining parts
are framed with the tables they will fill.

\subsection{Experimental Setup}
\label{sec:exp:setup}

\textbf{Datasets.} \corpus is described in \cref{tab:corpus}. No standard
benchmark exists for term type induction, so \corpus is both the input and, via
its citations, the ground against which faithfulness is checked. For RQ3 the
content is DLT-Corpus~\cite{hernandez2026dltcorpus}: from its first shard of
scientific papers (37,440 in the collection) and of USPTO patents (49,023) we
keep English items whose title or abstract names a DLT term, which 40\% of
papers and 49\% of patents do, and sample disjoint batches of 150 papers and
150 patents; an item is the title and abstract capped at 450 words. \todo{Held-out set and expert gold set for RQ4 and RQ6; contamination note.}

\textbf{Baselines.} \todo{Expert taxonomies; zero-shot LLM without the corpus; chunked RAG and
RAPTOR over it; TnT-LLM and LLMs4OL-style term typing; same corpus, model and
schema for all.}

\textbf{Metrics.} For faithfulness we report the share of citations whose quote
the verifier finds in the cited document (exact or relabelled), the share
relabelled, and the rejection rate; smaller rejection is better. For coverage we
report the number of term types, values and aliases, the number of documents
grounding each term type, and the number of term types with a recorded
conflict. For the hierarchy we report edges, roots, leaves, depth, bits per descent,
expected depth, AUC and effective leaves as defined in
\cref{sec:method:hierarchy}. \todo{For RQ4 add precision and recall of term
types against the expert gold set, and coverage of each expert taxonomy's
dimensions; for RQ6 add macro-F$_1$ of the downstream classifier.}

\textbf{Implementation.} Both tasks use \texttt{claude-haiku-5-5} called
through the Anthropic Messages API at effort \texttt{high} with adaptive
thinking disabled; the corpus sits behind a one-hour prompt-cache breakpoint.
Task~A is one streamed call with an output limit of 120,000 tokens; Task~B is
93 calls with $k{=}6$, four in parallel, each limited to 4,000 output tokens.
Task~C runs at effort \texttt{medium} with six items in flight; the discovery
and consolidation calls use the defaults of \cref{sec:method:maintain}. The verifier uses 12-character
shingles and an 80\% threshold. All runs are API calls from a laptop; no GPU is
involved. We report a single run. \todo{Add a second run with the same prompt
and report agreement of term types and values, and a run with
\texttt{claude-sonnet-5-5}, before submission.}

\subsection{RQ1: Term Types Generated and Their Faithfulness}
\label{sec:exp:rq1}

\begin{table}[t]
  \centering
  \caption{Term types by category. Docs is the mean number of source documents
  grounding a term type; Single counts term types grounded in one document;
  Confl.\ counts term types whose notes record a disagreement between sources;
  Verif.\ is the share of citations whose quote the verifier found in the cited
  document, exact or relabelled. Single run.}
  \label{tab:taska}
  \scriptsize
  %% Generated by paper/scripts/task_a_tables.py -- do not edit by hand.
\setlength{\tabcolsep}{3.5pt}
\begin{tabular}{@{}lrrrrrrr@{}}
    \toprule
    Category & Types & Values & Docs & Single & Confl. & Cites & Verif.\,\% \\
    \midrule
    Asset & 21 & 68 & 4.5 & 7 & 21 & 101 & 100.0 \\
    Technology & 24 & 82 & 3.2 & 5 & 16 & 123 & 97.6 \\
    Actor & 19 & 67 & 4.1 & 0 & 12 & 103 & 95.1 \\
    Activity & 10 & 43 & 3.9 & 1 & 8 & 63 & 100.0 \\
    Obligation or risk & 19 & 57 & 2.7 & 7 & 12 & 90 & 97.8 \\
    \midrule
    All & 93 & 317 & 3.6 & 20 & 69 & 480 & 97.9 \\
    \bottomrule
\end{tabular}

\end{table}


\textbf{What the corpus supports.} The run produced 93 term types with 317
values and 212 aliases (\cref{tab:taska}). We make the
following observations. (1) Technology dimensions are the largest group,
because the blockchain taxonomy of Tasca and Tessone contributes architectural
dimensions that no legal instrument uses, such as the fee regime and the
interoperability of a ledger. (2) Asset and actor dimensions are the most
widely shared. An asset term type is grounded in 4.5 documents on average, and
the status of an actor as a service provider is grounded in 13 of the 24
documents. (3) Twenty term types rest on a single document, mostly distinctions the
academic taxonomies draw and the law does not yet make, such as the mechanism
by which yield reaches a holder. Every document is cited by at least one term
type; the FATF guidance of 2021 is cited most, by 27, and the DLT-Corpus paper
least, by one; its sections describe data sources and metrics rather than
classification criteria.


Sixty-nine term types carry a note recording a disagreement between sources.
\Cref{tab:case} (\cref{app:results}) shows the reference-asset dimension of
\cref{fig:one} as returned: five values, seven grounding documents across four
jurisdictions, and a note recording the jurisdictional split. The
fungibility dimension likewise records that fungibility is an element of the UK
and US definitions of the asset class but not of the \src{FATF} or \src{MiCA}
definitions. These notes are the raw material for the jurisdiction-indexed
nodes of \cref{sec:method:hierarchy}.


\textbf{Faithfulness.} Of the 480 citations, 463 matched the cited section
verbatim under normalisation, 7 were relabelled, and 10 were rejected, a
rejection rate of 2.1\% (\cref{tab:taska}, last column).1\% for actor dimensions to 100\% for asset and
activity dimensions. One cited label did not exist, a subsection heading inside
the FinCEN guidance that the model used in place of the enclosing section's
label, and six named a real but wrong section of the right document. The verifier located the quote in each of these
seven cases. \Cref{tab:errors} in \cref{app:results} lists the ten rejections with a
diagnosis. Two
quotes share no five-word sequence with the cited section. Three compress a
definition, two drop the enumerated items of a statutory list, as in
``'blockchain' means any technology (i) where data is shared across a
network'', and three reproduce most of a sentence and then rephrase its end.
These quotes point at the right passage but are not verbatim, so we keep them in
the output as unverified rather than repairing them. Every term type keeps at
least one verified citation, but five values were supported only by a rejected
quote and are therefore unverified.

Verification is a necessary check, not a sufficient one. A verbatim quote can
still fail to support the dimension it is attached to, and our verifier does not
read for relevance. Whether the 470 verified citations do support their term
types is the question the expert annotation planned in \cref{sec:exp:setup} is
designed to answer, and we have not run it yet.

\subsection{RQ2: The Is-a Hierarchy}
\label{sec:exp:rq2}

\begin{table}[t]
  \centering
  \caption{The is-a DAG by category after consolidation, no weight threshold.
  Bits/desc.\ is the information gain of a random root-to-leaf descent;
  E[depth] its expected length; AUC the area under the reach-a-leaf curve (1 when
  every root is a leaf); Eff.\,lv.\ is the effective number of leaves
  $2^{\entropy(r)}$ for the leaf-reach distribution $r$. Single run.}
  \label{tab:taskb}
  \scriptsize
  %% Generated by paper/scripts/task_a_tables.py -- do not edit by hand.
\setlength{\tabcolsep}{2pt}
\begin{tabular}{@{}lrrrrrrrrr@{}}
    \toprule
    Category & Nodes & Edges & Roots & Leaves & Depth & Bits/desc. & E[depth] & AUC & Eff.\,lv. \\
    \midrule
    Asset & 21 & 68 & 1 & 5 & 9 & 4.17 & 2.54 & 0.773 & 3.2 \\
    Technology & 24 & 75 & 1 & 4 & 14 & 5.47 & 3.94 & 0.754 & 2.6 \\
    Actor & 19 & 58 & 1 & 3 & 12 & 6.02 & 3.30 & 0.767 & 2.6 \\
    Activity & 10 & 31 & 1 & 2 & 7 & 3.88 & 2.32 & 0.741 & 2.0 \\
    Obligation or risk & 19 & 47 & 1 & 4 & 12 & 4.05 & 1.93 & 0.881 & 4.0 \\
    \midrule
    All & 93 & 279 & 5 & 18 & 14 & 4.72 & 2.81 & 0.835 & 14.0 \\
    \bottomrule
\end{tabular}

\end{table}


\textbf{Structure.} The 93 calls returned \BWeighted\ weighted candidate edges,
of which consolidation kept \BEdges\ and dropped \BDropped\ for closing a cycle,
with \BTies\ exact tie (\cref{tab:taskb}). We make the following observations.
(1) Every category closes under a single root, the crypto-asset definitional
criteria for assets, the governance model for technologies, service-provider
status for actors, the service type for activities, and AML applicability for
obligations; the model was free to return \emph{none} for any node and did so
with median mass 0.13. (2) The DAG is deep and multi-parented: a node keeps
\BMeanParents\ parents on average and up to \BMaxParents, the technology category
reaches depth 14, and a random descent resolves \BBits\ bits before ending at
one of \BEffLeaves\ effective leaves. (3) In the asset category (\cref{fig:tree}, \cref{app:results}) the reference-asset dimension
specialises the legal classification and carries the stablecoin class,
redemption mechanism and tokenisation form beneath it, but second-choice
parents carry nearly as much mass: the model gave
\emph{legal classification} and \emph{stablecoin regulatory class} the same
weight of 0.32 as parents of \emph{reference asset}, and \BMutual\ pairs of
nodes were each proposed as the other's parent. The weights are coarse, with a
median top weight of \BMedianTop, so we read them as rankings rather than
calibrated probabilities.


\begin{figure*}[t]
  \centering
  %% Figure: hierarchy metrics when edges below a weight are hidden (text width, 1x3).
%% Legends sit above each panel so they never cover the curves.
\begin{tikzpicture}
  \begin{groupplot}[
    group style={group size=3 by 1, horizontal sep=15mm},
    width=0.33\textwidth, height=3.6cm,
    xmin=0, xmax=1, xtick={0,0.2,0.4,0.6,0.8,1},
    tick label style={font=\scriptsize}, label style={font=\scriptsize},
    ylabel near ticks, xlabel near ticks,
    axis line style={gray!60}, ymajorgrids, grid style={gray!25},
    legend style={font=\scriptsize, draw=none, fill=none, at={(0.5,1.02)}, anchor=south, legend columns=-1,
                  /tikz/every even column/.append style={column sep=5pt}},
    every axis plot/.append style={mark size=1.4pt, line width=0.9pt}]
    \nextgroupplot[ylabel={count}, xlabel={(a) minimum edge weight}, ymin=0]
      \addplot[blue!60!black, mark=*, mark repeat=10] table[x=thr, y=edges] {tables/spectrum.dat};
      \addplot[orange!80!black, mark=square*, mark repeat=10] table[x=thr, y=roots] {tables/spectrum.dat};
      \addplot[green!50!black, mark=triangle*, mark repeat=10] table[x=thr, y=leaves] {tables/spectrum.dat};
      \legend{kept edges, roots, leaves}
    \nextgroupplot[ylabel={bits per descent}, xlabel={(b) minimum edge weight}, ymin=0]
      \addplot[blue!60!black, mark=*, mark repeat=10] table[x=thr, y=bits] {tables/spectrum.dat};
    \nextgroupplot[ylabel={fraction}, xlabel={(c) minimum edge weight}, ymin=0, ymax=1.05, ytick={0,0.25,0.5,0.75,1}]
      \addplot[blue!60!black, mark=*, mark repeat=10] table[x=thr, y=auc] {tables/spectrum.dat};
      \addplot[orange!80!black, mark=square*, mark repeat=10] table[x=thr, y expr=\thisrow{effleaves}/93] {tables/spectrum.dat};
      \legend{AUC, effective leaves / 93}
  \end{groupplot}
\end{tikzpicture}

  \caption{Sensitivity to the display threshold: the sub-DAG keeps only edges
  with weight at or above the x value. (a) kept edges, roots and leaves; (b)
  information gain per descent; (c) AUC and effective leaves as a fraction of
  the 93 nodes. Hiding edges below 0.4 keeps 64 edges and leaves 30 roots.}
  \label{fig:spectrum}
\end{figure*}

\textbf{Sensitivity to the threshold.} \Cref{fig:spectrum} sweeps a display
threshold over the consolidated DAG. Hiding edges below 0.2 halves the edge
count and roughly halves the bits per descent; at 0.4 a third of the nodes are
roots and a descent resolves half a bit; above 0.6 the graph is a forest of
near-singletons. The full DAG carries the most discrimination; a thresholded
one is easier to read, and the curves price each cut.

\textbf{Faithfulness of edges.} The entries carry \BCites\ citations, of which
\BVerified\ (\BVerifiedPct\%) verify, \BRelabelled\ after relabelling, and
\BRejected\ are rejected, a markedly higher rejection rate than in Task~A. Only
\BEntriesCited\ of the \BWeighted\ entries cite anything, which the prompt
allows for low-probability entries, and \BTopNoCite\ top entries cite nothing,
which it does not. Rejected quotes are mostly short phrases cited from
long sections, such as ``Wrapped tokens''. Verification again checks
existence, not that the quote evidences the relation.
\todo{Expert judgement of a sample of 100 edges on a 1--3 scale with agreement,
following the protocol planned for \cref{sec:exp:rq4}.}

\subsection{RQ3: The DAG as a Label Space}
\label{sec:exp:rq3}

\begin{table*}[t]
  \centering
  \caption{Two maintenance iterations. Each tags a disjoint batch with the
  catalogue version in force, measures the DAG on it, discovers candidates at
  the target nodes and consolidates the accepted ones into the next version.
  Evid.\ is the share of evidence spans found verbatim in the item; Nodes used
  counts nodes that received at least one item; USD and minutes cover tagging,
  discovery and consolidation at list price with eight workers.}
  \label{tab:evolution}
  \scriptsize
  %% Generated by paper/scripts/task_a_tables.py -- do not edit by hand.
\setlength{\tabcolsep}{3pt}
\begin{tabular}{@{}llrrrrrrlllr@{}}
    \toprule
    It. & Cat. & Items & Tags/item & Evid.\,\% & Nodes used & Collapsed & Targets & Cand.\ (acc.) & Produced: nodes/edges/values & USD & Min. \\
    \midrule
    1 & v0 & 300 & 3.55 & 96.3 & 80/93 & 3 & 15 & 32 (8) & v1: 99/285/351 & 0.15 & 2 \\
    2 & v1 & 8,646 & 4.06 & 96.8 & 98/99 & 11 & 32 & 81 (21) & v2: 107/293/396 & 4.29 & 34 \\
    \bottomrule
\end{tabular}

\end{table*}

\begin{table}[t]
  \centering
  \caption{Objectives of each catalogue version from its is-a weights (prior)
  and from the tags of the batch classified with it.}
  \label{tab:objectives}
  \scriptsize
  %% Generated by paper/scripts/task_a_tables.py -- do not edit by hand.
\setlength{\tabcolsep}{3pt}
\begin{tabular}{@{}lrrrrrrrrrrrr@{}}
    \toprule
    Objective & v0 prior & iter.\ 1 tags & v1 prior & iter.\ 2 tags & v2 prior & iter.\ 3 tags & v3 prior & iter.\ 4 tags & v4 prior & iter.\ 5 tags & v5 prior & iter.\ 6 tags \\
    \midrule
    Bits per descent & 4.72 & 6.92 & 4.82 & 6.96 & 4.84 & 7.02 & 4.88 & 7.22 & 5.01 & 7.27 & 5.02 & 7.06 \\
    AUC of reach-a-leaf & 0.835 & 0.794 & 0.843 & 0.771 & 0.851 & 0.778 & 0.843 & 0.780 & 0.833 & 0.773 & 0.833 & 0.773 \\
    Expected depth & 2.81 & 3.37 & 2.70 & 3.70 & 2.74 & 3.83 & 2.86 & 3.80 & 3.00 & 3.90 & 3.00 & 3.90 \\
    Effective leaves & 14.0 & 19.7 & 18.5 & 21.3 & 22.5 & 21.7 & 23.1 & 24.7 & 25.7 & 24.8 & 26.1 & 22.9 \\
    Effective leaves / nodes & 0.150 & 0.212 & 0.187 & 0.215 & 0.210 & 0.203 & 0.204 & 0.218 & 0.207 & 0.200 & 0.204 & 0.179 \\
    \bottomrule
\end{tabular}

\end{table}

\textbf{Tagging.} \Cref{tab:evolution} summarises both iterations. The pilot
batch of \CItems\ items received \CTags\ tags, \CTagsPerItem\ per item, and the
second batch of \TItems\ items, a tenth of the on-topic collection, received
\TTags\ tags, \TTagsPerItem\ per item; \CEvidencePct\% and \TEvidencePct\% of
the evidence spans were found verbatim in the item, and \CRejected\ and
\TRejected\ tags were rejected for unknown identifiers. We make the following
observations. (1) The content exercises the technical half of the DAG: in the
second batch \TCovTech\ items carry a technology tag and \TCovAsset\ an asset
tag, against \TCovActor\ for actors, \TCovActivity\ for activities and
\TCovOblig\ for obligations. The legal nodes are dormant on the pilot batch and collapsed on a single
value on the larger one; testing those branches needs regulatory content. (2) The dimensions discovered in
iteration 1 are among the most used in iteration 2: ledger purpose was tagged
on \TLedgerPurpose\ items with a value entropy of \TEntLedgerPurpose\ bits,
smart-contract usage domain on \TUsageDomain, identity binding on \TIdentity\
and the privacy-mechanism family on \TPrivacy\ with \TEntPrivacy\ bits, so
they discriminate the content rather than absorbing it. (3) Tagging cost
\$\TCostPerItem\ per item at \TMeanLat\ seconds per request; the whole
on-topic collection of about \OnTopicK\ thousand documents would cost about
\$\ProjCost\ and \ProjHours\ hours at the same concurrency.

\textbf{Measurement.} \Cref{tab:objectives} compares, for each version, the
objectives from the model's is-a weights with those measured from the tags. A
content-driven descent resolves about seven bits against under five from the
prior and ends at more effective leaves: the content goes deeper and spreads
wider than the prior weights suggested. Collapsed nodes grow with the batch,
\TCollapsed\ in iteration 2, led by the DLT data structure with 96\% of 2,377
items on ``blockchain''; they are also the largest residual pools, which is
where discovery is pointed.

\begin{table}[t]
  \centering
  \caption{Term types added by the two iterations, with the node they were
  discovered under, the number of shown residual items supporting them and the
  increase in branching gain at that node in bits. Values added to existing
  nodes (\EAddedValues\ and \TAddedValues) are not listed.}
  \label{tab:accepted}
  \scriptsize
  %% Generated by paper/scripts/task_a_tables.py -- do not edit by hand.
\setlength{\tabcolsep}{3pt}
\begin{tabularx}{\columnwidth}{@{}lX>{\raggedright\arraybackslash}p{0.3\columnwidth}rr@{}}
    \toprule
    It. & New term type & Under & Items & Gain \\
    \midrule
    1 & Privacy mechanism family & Privacy and anonymity features & 8 & 0.83 \\
    1 & Smart contract deployment context & Smart contract / script language expressiveness & 6 & 0.66 \\
    1 & Smart contract usage domain & Smart contract / distributed ledger application scope & 8 & 0.46 \\
    1 & Identity binding mechanism & Identity layer and onboarding & 5 & 0.34 \\
    1 & Ledger purpose & Recording and cryptographic linkage & 6 & 0.21 \\
    1 & Ledger partitioning or layering topology & DLT data structure (blockchain vs non-chain DLT) & 5 & 0.18 \\
    2 & Cross-chain transfer mechanism & Wrapped / cross-chain representation & 8 & 0.74 \\
    2 & Key custody and control model & Wallet storage mode (hot or cold) & 7 & 0.67 \\
    2 & Finality mechanism & Finality of agreement & 5 & 0.55 \\
    2 & Key isolation from network connectivity & Hosted versus unhosted (self-hosted) wallet & 5 & 0.55 \\
    2 & Payment settlement mechanism & Payment service and funds transfer & 5 & 0.38 \\
    2 & Key splitting scheme & Permitted control over client assets (unilateral vs multi-party control) & 8 & 0.32 \\
    2 & Expressiveness claim basis & Smart contract / script language expressiveness & 5 & 0.25 \\
    2 & Key management control model & Crypto-asset service type & 6 & 0.11 \\
    3 & Zero-knowledge proof system variant & Privacy mechanism family & 6 & 0.61 \\
    3 & Key share distribution topology & Key management control model & 5 & 0.59 \\
    3 & Fiat-rail integration & Payment settlement mechanism & 5 & 0.55 \\
    3 & Key recovery mechanism & Hosted versus unhosted (self-hosted) wallet & 6 & 0.50 \\
    3 & Script verification approach & Smart contract / script language expressiveness & 5 & 0.26 \\
    3 & Backing of token & Tokenisation form & 5 & 0.14 \\
    4 & Deanonymisation / tracing approach & Anonymity-enhancing crypto-asset status & 6 & 0.75 \\
    4 & Backing mechanism & Backing of token & 5 & 0.72 \\
    4 & Settlement finality timing & Payment settlement mechanism & 5 & 0.42 \\
    4 & Settlement counterparty & Fiat-rail integration & 5 & 0.65 \\
    4 & Physical transfer mechanism across air gap & Key isolation from network connectivity & 5 & 0.57 \\
    4 & Validator selection mechanism & Validator operation capacity (who may run a node) & 5 & 0.55 \\
    4 & Sharding partition basis & Ledger partitioning or layering topology & 5 & 0.54 \\
    4 & Wallet interface or delivery form & Hosted versus unhosted (self-hosted) wallet & 6 & 0.44 \\
    4 & Share recovery mechanism & Key splitting scheme & 5 & 0.44 \\
    4 & Centralisation locus & Centralisation of control & 6 & 0.38 \\
    4 & Cross-chain interoperability mechanism & Wrapped / cross-chain representation & 7 & 0.28 \\
    5 & Key isolation mechanism & Key isolation from network connectivity & 7 & 0.59 \\
    5 & Key distribution method across isolation boundary & Key isolation from network connectivity & 5 & 0.48 \\
    5 & Finality guarantee mode & Finality of agreement & 5 & 0.41 \\
    5 & Key share cryptographic primitive & Key splitting scheme & 5 & 0.30 \\
    6 & ZKP application role & Zero-knowledge proof system variant & 5 & 0.86 \\
    6 & Key storage medium or environment & Key custody and control model & 5 & 0.79 \\
    6 & Parent-child coupling mechanism & Ledger partitioning or layering topology & 6 & 0.58 \\
    6 & Finality guarantee source & Finality of agreement & 5 & 0.47 \\
    \bottomrule
\end{tabularx}

\end{table}

\textbf{Discovery and evolution.} Iteration 1 found \MTargets\ targets and
accepted \MAccepted\ of \MCandidates\ candidates; iteration 2 found \TTargets\
targets and accepted \TAccepted\ of \TCandidates, producing version 2 with
\VTwoNodes\ term types, \VTwoEdges\ edges and \VTwoValues\ values
(\cref{tab:accepted}). The additions are dimensions a legal corpus would not
contain, such as the cryptographic family behind a privacy guarantee or the
custody model of a key, and several iteration-2 values refine iteration-1
nodes. Two failures are visible. Two application-domain dimensions proposed
under two smart-contract nodes in iteration 1 ask the same question and were
not merged by the consolidation call. And with raw support counts, capped by
the forty items shown, only 10 of 76 candidates cleared the gain test against
residual pools of thousands in iteration 2; extrapolating support to the
residual is the rule we report, and the rerun accepted 21 of 81.

The tags have no gold labels yet, so verified evidence is the only quality
signal we report; an evidence span can exist without the tag being right.
\todo{Expert review of a slice of the tags for precision and recall, and a
regulatory content batch for the legal branches.}

\subsection{Planned Evaluations (RQ4--RQ6)}
\label{sec:exp:rq4}
\label{sec:exp:rq5}
\label{sec:exp:rq6}
\todo{RQ4: expert taxonomies and LLM baselines (zero-shot, chunked RAG, RAPTOR,
TnT-LLM) on the same corpus; coverage, new dimensions accepted by experts,
verification rate. RQ5: ablation rows for grounding strategy, model and prompt
variants beside the manifest result of \cref{sec:exp:rq1}. RQ6: macro-F$_1$ of
classifying Web-sourced token descriptions with and without the ontology, with
one failure case.}

\subsection{Cost}
\label{sec:exp:cost}
Task~A is one call of 302 seconds with 85,729 output tokens, 2,542 uncached
input tokens and 235,298 tokens read from cache: at most \$0.33 at list price
with cached tokens charged in full. Task~B is \BCalls\ calls in \BSeconds\
seconds with \BOutTokens\ output tokens and \BCacheTokens\ million cached
tokens read, at most \$11.7 before the cache-read discount. \todo{Exact billing figures; baseline and frontier-model cost.}

